% sections/literature.tex: Section 2: what each literature finds and what this paper adds.
\section{Related Literature}\label{sec:literature}

\paragraph{Discount rates from asset markets.} The standard approach to measuring time preferences uses the Euler equation to link the price of risk-free bonds to the discount factor of the marginal investor \citep{lustig_market_2005, bansal_risks_2004}. Calibrations of heterogeneous-agent models set $\delta \approx 0.96$ to match the capital-output ratio \citep{guvenen_reconciling_2006} or government bond yields, and \citet{di_tella_zero-beta_2024} construct a zero-beta rate from equity portfolios and estimate an average of 8.3 percent per year. These approaches identify the preferences of the marginal investor, presumably a wealthy and patient household, which may differ substantially from those of the typical household. This paper identifies the discount rate of the borrowing population, the majority of U.S.\ households, and finds it near the zero-beta rate and well above the calibrated factor.

\paragraph{Laboratory and field elicitations.} \citet{coller_eliciting_1999} elicit discount rates from students choosing between dated payments and find rates near 20 percent once the field interest rate is provided as a reference; \citet{harrison_estimating_2002} elicit 28 percent from a representative Danish sample over horizons of six months to three years. \citet{cohen_measuring_2020} survey 222 studies through 2016; the 112 that report a point estimate have a median of 22 percent per year and a mean of 26 percent, primarily from small samples with small stakes, and behavior under small stakes may differ from behavior under large stakes \citep{rabin_risk_2000}. The elicitations are over horizons of weeks to a few years, where a payment due now already differs from a payment due later by the premium this paper measures; the exponential rate this paper estimates over the contract horizon is a fraction of the elicited rates, and Section \ref{sec:comparison} states how the two segments of the schedule relate.

\paragraph{Default, payment relief and household debt.} \citet{ganong_liquidity_2020} use two regression discontinuities in the HAMP program and find that payment reductions prevent mortgage default far more effectively than principal reductions of equal cost. \citet{dobbie_targeted_2020} randomize the terms of debt-management plans and find that a write-down of the last payments lowers bankruptcy while an immediate payment reduction that adds payments later raises it. \citet{indarte_moral_2023} separates the cash-on-hand and moral-hazard channels of bankruptcy with adjustable-rate resets and exemption kinks and finds the cash-on-hand response an order of magnitude larger per dollar. Each of these designs is a payment path, and Section \ref{sec:quasi} applies the formula of this paper to their pairs: the first two identify a discount rate for the populations they treat, the third the marginal-utility ratio. \citet{ganong_consumer_2019} fit consumption during unemployment and find that about one third of consumers behave as impatient low-$\beta$ types, without a distribution of rates; this paper supplies the rate by group at population scale. The option-value view of default, in which the borrower who continues keeps the right to default later \citep{deng_mortgage_2000, campbell_model_2015}, is built into the model of Section \ref{sec:model} and is what prevents the decline of the payment response over the life of a loan from being a measure of the rate (Appendix \ref{appendix:mcliquidity}). \citet{athreya_persistence_2019} reproduce the persistence of financial distress within the person only with persistent heterogeneity in discount factors; this paper measures the person's share of the variance of default at a quarter and finds no gradient in the rate across the groups whose default rates differ most (Appendix \ref{appendix:traits}, Section \ref{sec:rate_people}).

\paragraph{Sufficient statistics.} The sufficient-statistics approach to welfare and policy analysis \citep{chetty_sufficient_2009} recovers policy-relevant parameters from a small number of estimable moments without specifying the full structure of the model. \citet{davila_exemptions_2020} applies it to household default: the optimal bankruptcy exemption depends on four measurable objects, among them the value households place on a marginal dollar in the bankruptcy state, which he recovers from consumption changes under an assumed utility function and an assumed annual discount factor of $0.94$. The discount function that paper assumes is the object this paper measures. The same ratio logic identifies a marginal-utility ratio at one date in \citet{chetty_moral_2008} and \citet{indarte_moral_2023}; this paper's increment is the ratio across dates, and Appendix \ref{appendix:priority} states the antecedent and the increment for each result.

\paragraph{Discount rates in public finance.} Governments publish declining discount-rate schedules for long horizons \citep{arrow_should_2014}, the U.S.\ government recently lowered its recommended social discount rate from 3 to 2 percent \citep{the_white_house_biden-harris_2023}, and \citet{hendren_unified_2020} show that the ranking of public investments turns on the rate, with adult programs dominating at high rates and early-childhood programs at low rates; \citet{demarzo_sovereign_2023} show that patient citizens are harmed when their government borrows excessively. This paper places the contract-horizon discount rate of borrowing households at $0.11$, above the two percent rate used in policy analyses and inside the set the rule-based experiments accept for distressed borrowers, with a further premium on the payment due now.
